\section{Discussion}
\label{sec:discussion}

The results support partner-local affective precision as a calibration mechanism for social policy commitment. Precision was defined from partner-response predictability rather than payoff, lost its policy-sharpening effect when the $\beta_k \rightarrow \gamma_k$ pathway was cut, and remained behaviorally active under abrupt change without guaranteeing payoff improvement. Together, these findings support the paper's central distinction: social prediction and social commitment are related but distinct. A reliably cooperative partner and a reliably defecting partner can both generate low partner-response likelihood surprisal because both are predictable; conversely, a favorable but inconsistent partner can remain difficult to model. Partner-local affective precision therefore estimates how well the current partner model is working, then regulates how strongly that model is expressed in policy.

The tracked-only ablation clarifies where the mechanism acts. Updating $\beta_k$ without allowing it to modulate $\gamma_k$ preserved the tracker but removed the same policy-sharpening effect, while full precision modulation sharpened policy commitment. This supports the interpretation that partner-local affective precision does not change the partner-state update rule directly; any accuracy differences that emerge under partner choice are downstream sampling effects, reflecting which partners the agent engages and observes, while the mechanism itself changes how strongly current partner beliefs are deployed. It also clarifies the distinction from affective-precision accounts centered on policy-inference quantities: here, the model-fitness signal is partner-response likelihood surprisal under the pre-update predictive type--stance posterior \(q_{k,t}^-\).

The abrupt-betrayal condition shows an important limitation of confidence accumulated from past predictive success. Confidence built from reliable evidence can remain behaviorally active after the relationship changes, producing a lag between changes in partner behavior and changes in policy commitment. Future extensions could incorporate volatility or change detection that discounts accumulated confidence when prediction-error structure shifts, or adapt the confidence baseline itself to a partner's recent predictive history. The present implementation maintains $\beta_k$ as an auxiliary tracker; representing confidence within a deeper generative hierarchy is another possible extension, but is not required for the relationship-specific calibration mechanism studied here.

The profile analyses generate computational hypotheses about trust calibration, with potential links to work on social learning, psychosis, and attachment~\cite{Adams2013ComputationalPsychosis,Behrens2008SocialLearning,Cittern2018AttachmentActiveInference}. Precision gain and prior model fitness shaped update amplitude, sensitivity to change, and reengagement after repair, but no single profile was uniformly best across the reported readouts. Higher gain expanded confidence swings without monotonic payoff improvement, while different gain--prior combinations produced distinct patterns of confidence accumulation, betrayal adaptation, and repair. These tradeoffs suggest empirical tests on trust-game time series, including repair settings after trust violations: whether human behavior separates predictability from payoff, policy deployment from partner-state inference, and reengagement from restored confidence. A further extension is to replace scripted partners with reciprocal active-inference agents that maintain their own beliefs, policies, and precision estimates.

\section{Conclusion}
\label{sec:conclusion}
\enlargethispage{3\baselineskip}

We formalize social affective precision as a relationship-specific metacognitive signal: an estimate of how confidently an agent should select policies from its current model of each partner. Implemented as partner-local affective precision, this signal updates from partner-response predictive evidence and modulates policy precision during selection. Across simulations in a graded multi-partner trust game, precision tracks predictability more closely than payoff, affects behavior through the $\beta_k \rightarrow \gamma_k$ policy-precision pathway rather than by directly modifying partner-state inference, and remains behaviorally active when relationships change.

These results support a distinction between social prediction and social commitment: an agent can infer what a partner is likely to do while separately estimating how strongly that inference should guide policy commitment. Precision gain and prior model fitness shape this calibration process, producing tradeoffs in confidence accumulation, betrayal adaptation, reengagement, and repair. Future empirical work can test whether human trust-game behavior similarly separates predictability from payoff, inference from policy commitment, and reengagement from restored confidence.
